\chapter{Index des méthodes}

\noindent\textit{Les 68 méthodes du fascicule, regroupées par chapitre, pour retrouver
rapidement « comment faire » sans connaître le numéro du chapitre.}
\vspace{8pt}

\begin{multicols}{2}
\footnotesize

\textbf{\normalsize\color{onbo}Chapitre 1 --- Les constituants de la matière}
\begin{enumerate}[label=\arabic*.,itemsep=1pt,topsep=2pt,leftmargin=1.6em]
\item Écrire le symbole d'un atome usuel
\item Compter les atomes d'une molécule à partir de sa formule brute
\item Distinguer corps simple et corps composé
\item Écrire la formule brute d'une molécule connaissant sa composition
\item Déterminer le symbole d'un ion à partir d'un atome et du nombre d'électrons
\item Reconnaître une liaison de covalence dans une molécule
\end{enumerate}
\vspace{4pt}

\textbf{\normalsize\color{onbo}Chapitre 2 --- Classification périodique des éléments}
\begin{enumerate}[label=\arabic*.,itemsep=1pt,topsep=2pt,leftmargin=1.6em]
\item Déterminer le nombre d'électrons d'un atome neutre à partir de son numéro
\item Comparer deux éléments à l'aide de leur numéro atomique
\item Identifier la famille chimique d'un élément
\item Prévoir une propriété par comparaison au sein d'une famille
\item Reconnaître un gaz noble
\end{enumerate}
\vspace{4pt}

\textbf{\normalsize\color{onbo}Chapitre 3 --- La mole et quantité de matière}
\begin{enumerate}[label=\arabic*.,itemsep=1pt,topsep=2pt,leftmargin=1.6em]
\item Calculer une masse molaire moléculaire
\item Calculer une quantité de matière à partir d'une masse
\item Calculer une masse à partir d'une quantité de matière
\item Calculer un nombre d'entités à partir d'une quantité de matière
\item Calculer une quantité de matière à partir d'un nombre d'entités
\item Résoudre un problème concret combinant masse, quantité de matière et nombre
\end{enumerate}
\vspace{4pt}

\textbf{\normalsize\color{onbo}Chapitre 4 --- Réactions chimiques et loi de Lavoisier}
\begin{enumerate}[label=\arabic*.,itemsep=1pt,topsep=2pt,leftmargin=1.6em]
\item Identifier réactifs et produits dans un énoncé
\item Appliquer la loi de Lavoisier pour calculer une masse manquante
\item Vérifier qu'une équation-bilan est équilibrée
\item Équilibrer une équation-bilan simple
\item Résoudre un problème concret avec la loi de Lavoisier
\end{enumerate}
\vspace{4pt}

\textbf{\normalsize\color{onbo}Chapitre 5 --- Électrolyse et synthèse de l'eau}
\begin{enumerate}[label=\arabic*.,itemsep=1pt,topsep=2pt,leftmargin=1.6em]
\item Écrire l'équation-bilan de l'électrolyse de l'eau
\item Calculer le volume de dioxygène connaissant le volume de dihydrogène
\item Calculer le volume de dihydrogène connaissant le volume de dioxygène
\item Identifier un gaz recueilli lors de l'électrolyse
\item Écrire et vérifier l'équation-bilan de la synthèse de l'eau
\end{enumerate}
\vspace{4pt}

\textbf{\normalsize\color{onbo}Chapitre 6 --- Les solutions aqueuses}
\begin{enumerate}[label=\arabic*.,itemsep=1pt,topsep=2pt,leftmargin=1.6em]
\item Distinguer soluté et solvant dans une solution
\item Déterminer si une solution est conductrice
\item Calculer une concentration molaire
\item Déterminer la nature d'une solution à partir de son pH
\item Identifier un ion par un test caractéristique
\end{enumerate}
\vspace{4pt}

\textbf{\normalsize\color{onbo}Chapitre 7 --- Machines simples}
\begin{enumerate}[label=\arabic*.,itemsep=1pt,topsep=2pt,leftmargin=1.6em]
\item Calculer le poids d'un objet
\item Calculer la force sur un plan incliné
\item Calculer la force avec une poulie ou un palan
\item Calculer la force avec un treuil
\item Comparer l'intérêt de deux machines simples pour une même tâche
\end{enumerate}
\vspace{4pt}

\textbf{\normalsize\color{onbo}Chapitre 8 --- Loi d'Ohm et circuits électriques}
\begin{enumerate}[label=\arabic*.,itemsep=1pt,topsep=2pt,leftmargin=1.6em]
\item Appliquer la loi d'Ohm
\item Étudier un circuit en série
\item Étudier un circuit en dérivation
\item Calculer une puissance électrique
\end{enumerate}
\vspace{4pt}

\textbf{\normalsize\color{onbo}Chapitre 9 --- Production et transport de l'énergie électrique}
\begin{enumerate}[label=\arabic*.,itemsep=1pt,topsep=2pt,leftmargin=1.6em]
\item Calculer les pertes en ligne
\item Calculer la tension au secondaire d'un transformateur
\item Reconnaître un transformateur élévateur ou abaisseur
\item Utiliser la conservation de puissance d'un transformateur idéal
\end{enumerate}
\vspace{4pt}

\textbf{\normalsize\color{onbo}Chapitre 10 --- Tension alternative et énergie électrique domestique}
\begin{enumerate}[label=\arabic*.,itemsep=1pt,topsep=2pt,leftmargin=1.6em]
\item Calculer une période ou une fréquence
\item Calculer une énergie électrique
\item Calculer le coût d'une consommation électrique
\end{enumerate}
\vspace{4pt}

\textbf{\normalsize\color{onbo}Chapitre 11 --- Le pétrole et le gaz naturel}
\begin{enumerate}[label=\arabic*.,itemsep=1pt,topsep=2pt,leftmargin=1.6em]
\item Ordonner les fractions du pétrole selon leur volatilité
\item Écrire et vérifier une équation de combustion complète
\item Distinguer combustion complète et incomplète
\end{enumerate}
\vspace{4pt}

\textbf{\normalsize\color{onbo}Chapitre 12 --- Les matières plastiques}
\begin{enumerate}[label=\arabic*.,itemsep=1pt,topsep=2pt,leftmargin=1.6em]
\item Distinguer monomère et polymère
\item Distinguer thermoplastique et thermodurcissable
\item Analyser l'impact environnemental d'un objet plastique
\end{enumerate}
\vspace{4pt}

\textbf{\normalsize\color{onbo}Chapitre 13 --- Dessin technique}
\begin{enumerate}[label=\arabic*.,itemsep=1pt,topsep=2pt,leftmargin=1.6em]
\item Calculer une longueur sur le dessin à partir de l'échelle
\item Calculer une longueur réelle à partir du dessin
\item Identifier un type de trait
\item Reconnaître la disposition des trois vues normalisées
\end{enumerate}
\vspace{4pt}

\textbf{\normalsize\color{onbo}Chapitre 14 --- Transmission du mouvement de rotation}
\begin{enumerate}[label=\arabic*.,itemsep=1pt,topsep=2pt,leftmargin=1.6em]
\item Calculer la vitesse de sortie d'un engrenage
\item Calculer la vitesse de sortie d'un système poulies-courroie
\item Déterminer si un système est réducteur ou multiplicateur
\item Reconnaître le sens de rotation d'un système
\end{enumerate}
\vspace{4pt}

\textbf{\normalsize\color{onbo}Chapitre 15 --- Moteurs à combustion interne et moteurs électriques}
\begin{enumerate}[label=\arabic*.,itemsep=1pt,topsep=2pt,leftmargin=1.6em]
\item Ordonner les quatre temps du cycle moteur
\item Distinguer moteur essence et moteur diesel
\item Calculer un rendement
\end{enumerate}
\vspace{4pt}

\textbf{\normalsize\color{onbo}Chapitre 16 --- Installations électriques domestiques et maintenance}
\begin{enumerate}[label=\arabic*.,itemsep=1pt,topsep=2pt,leftmargin=1.6em]
\item Calculer la puissance maximale admissible sur un circuit
\item Vérifier si un ensemble d'appareils peut être branché sans risque
\item Distinguer les rôles des différents dispositifs de sécurité
\end{enumerate}
\vspace{4pt}

\end{multicols}
